% GENERATED FILE. Edit canonical lessons, then run npm run generate:books.
% canonical-source-hash: fnv1a64:55e4c1a80d1d9b39
% canonical-lessons: ES-C539-semestre, ES-C539-quincena, ES-C539-lustro, ES-C539-milenio, ES-C539-periodo

\chapter{Semestres y quincenas}
\label{ch:es-semestres-y-quincenas}
\hlchaptermodality{539}

\begin{chapteropening}
\textbf{By the end of this chapter:} \emph{I can talk about stretches of time: half a year, a fortnight, five years, a millennium and a period.}

Complete the last lesson of chapter 539: I can talk about stretches of time: half a year, a fortnight, five years, a millennium and a period.
\end{chapteropening}

\section[el semestre]{el semestre — the semester, the half-year}

\label{lesson:ES-C539-semestre}



\begin{quote}
Before the new one: say the Spanish for recently, then the Spanish for lately.
\end{quote}

\subsection*{You'll want to know: el semestre}
\textbf{el semestre} — \textquotedblleft{}the semester, the half-year\textquotedblright{}.

\begin{quote}
El primer semestre del año fue muy difícil.
\end{quote}

\begin{quote}
\emph{The first half of the year was very hard.}
\end{quote}

At university it is a term of six months.

\begin{cousinweb}
\emph{semestre} comes from Latin \emph{sēmestris} 'of six months' (EtymDB), built on \emph{sex} 'six' and \emph{mēnsis} 'month'.
\end{cousinweb}

\begin{grammarlens}[title={What you've built}]
One more word for telling, explaining and describing at B1.
\end{grammarlens}

\subsection*{Your turn}
\begin{itemize}
\raggedright
  \item \emph{Say it:} \emph{el semestre}
  \item \emph{Say it:} \emph{el semestre}, once more
  \item \emph{Say it:} a sentence of your own with \emph{el semestre}
  \item \emph{Recall:} say the Spanish for recently, then the Spanish for lately, then say \emph{el semestre} again
\end{itemize}

\begin{tcolorbox}[breakable,colback=teal!4,colframe=teal!35!black,title={Before you move on}]
What does el semestre mean? (\textquotedblleft{}The semester, the half-year\textquotedblright{}.) Say it once more.
\end{tcolorbox}
\section[la quincena]{la quincena — the fortnight, two weeks}

\label{lesson:ES-C539-quincena}



\begin{quote}
Before the new one: say the Spanish for lately, then the Spanish for the semester.
\end{quote}

\subsection*{You'll want to know: la quincena}
\textbf{la quincena} — \textquotedblleft{}the fortnight, two weeks\textquotedblright{}.

\begin{quote}
Nos fuimos de vacaciones la primera quincena de agosto.
\end{quote}

\begin{quote}
\emph{We went on holiday in the first two weeks of August.}
\end{quote}

Some workers are paid \emph{por quincenas}, every two weeks.

\begin{cousinweb}
\emph{quincena} is built on \emph{quince} 'fifteen' (EtymDB), which is inherited from Latin \emph{quīndecim} 'fifteen', from \emph{quīnque} 'five' and \emph{decem} 'ten' (Lewis \& Short). A fortnight is counted as fifteen days.
\end{cousinweb}

\begin{grammarlens}[title={What you've built}]
One more word for telling, explaining and describing at B1.
\end{grammarlens}

\subsection*{Your turn}
\begin{itemize}
\raggedright
  \item \emph{Say it:} \emph{la quincena}
  \item \emph{Say it:} \emph{la quincena}, once more
  \item \emph{Say it:} a sentence of your own with \emph{la quincena}
  \item \emph{Recall:} say the Spanish for lately, then the Spanish for the semester, then say \emph{la quincena} again
\end{itemize}

\begin{tcolorbox}[breakable,colback=teal!4,colframe=teal!35!black,title={Before you move on}]
What does la quincena mean? (\textquotedblleft{}The fortnight, two weeks\textquotedblright{}.) Say it once more.
\end{tcolorbox}
\section[el lustro]{el lustro — the five-year period}

\label{lesson:ES-C539-lustro}



\begin{quote}
Before the new one: say the Spanish for the semester, then the Spanish for the fortnight.
\end{quote}

\subsection*{You'll want to know: el lustro}
\textbf{el lustro} — \textquotedblleft{}the five-year period\textquotedblright{}.

\begin{quote}
En el último lustro la ciudad cambió mucho.
\end{quote}

\begin{quote}
\emph{In the last five years the city changed a lot.}
\end{quote}

A formal word, common in the press.

\begin{cousinweb}
\emph{lustro} is a learned borrowing from Latin \emph{lūstrum} (EtymDB), the cleansing sacrifice that Rome held every five years after the census.
\end{cousinweb}

\begin{grammarlens}[title={What you've built}]
One more word for telling, explaining and describing at B1.
\end{grammarlens}

\subsection*{Your turn}
\begin{itemize}
\raggedright
  \item \emph{Say it:} \emph{el lustro}
  \item \emph{Say it:} \emph{el lustro}, once more
  \item \emph{Say it:} a sentence of your own with \emph{el lustro}
  \item \emph{Recall:} say the Spanish for the semester, then the Spanish for the fortnight, then say \emph{el lustro} again
\end{itemize}

\begin{tcolorbox}[breakable,colback=teal!4,colframe=teal!35!black,title={Before you move on}]
What does el lustro mean? (\textquotedblleft{}The five-year period\textquotedblright{}.) Say it once more.
\end{tcolorbox}
\section[el milenio]{el milenio — the millennium}

\label{lesson:ES-C539-milenio}



\begin{quote}
Before the new one: say the Spanish for the fortnight, then the Spanish for the five-year period.
\end{quote}

\subsection*{You'll want to know: el milenio}
\textbf{el milenio} — \textquotedblleft{}the millennium\textquotedblright{}.

\begin{quote}
La ciudad tiene más de un milenio.
\end{quote}

\begin{quote}
\emph{The city is more than a thousand years old.}
\end{quote}

\begin{cousinweb}
\emph{milenio} is a learned borrowing from New Latin \emph{millennium}, built on \emph{mīlle} 'a thousand' and \emph{annus} 'year' (Lewis \& Short).
\end{cousinweb}

\begin{grammarlens}[title={What you've built}]
One more word for telling, explaining and describing at B1.
\end{grammarlens}

\subsection*{Your turn}
\begin{itemize}
\raggedright
  \item \emph{Say it:} \emph{el milenio}
  \item \emph{Say it:} \emph{el milenio}, once more
  \item \emph{Say it:} a sentence of your own with \emph{el milenio}
  \item \emph{Recall:} say the Spanish for the fortnight, then the Spanish for the five-year period, then say \emph{el milenio} again
\end{itemize}

\begin{tcolorbox}[breakable,colback=teal!4,colframe=teal!35!black,title={Before you move on}]
What does el milenio mean? (\textquotedblleft{}The millennium\textquotedblright{}.) Say it once more.
\end{tcolorbox}
\section[el período]{el período — the period}

\label{lesson:ES-C539-periodo}



\begin{quote}
Before the new one: say the Spanish for the five-year period, then the Spanish for the millennium.
\end{quote}

\subsection*{You'll want to know: el período}
\textbf{el período} — \textquotedblleft{}the period\textquotedblright{}.

\begin{quote}
Fue un período difícil para la familia.
\end{quote}

\begin{quote}
\emph{It was a difficult period for the family.}
\end{quote}

\emph{Periodo} without the accent is also correct.

\begin{cousinweb}
\emph{período} is a learned borrowing from Latin \emph{periodus}, which comes from Greek \emph{períodos} 'a going round, a cycle' (EtymDB), built on \emph{perí} 'around' and \emph{hodós} 'way'.
\end{cousinweb}

\begin{grammarlens}[title={What you've built}]
One more word for telling, explaining and describing at B1.
\end{grammarlens}

\subsection*{Your turn}
\begin{itemize}
\raggedright
  \item \emph{Say it:} \emph{el período}
  \item \emph{Say it:} \emph{el período}, once more
  \item \emph{Say it:} a sentence of your own with \emph{el período}
  \item \emph{Recall:} say the Spanish for the five-year period, then the Spanish for the millennium, then say \emph{el período} again
\end{itemize}

\begin{tcolorbox}[breakable,colback=teal!4,colframe=teal!35!black,title={Before you move on}]
What does el período mean? (\textquotedblleft{}The period\textquotedblright{}.) Say it once more.
\end{tcolorbox}
